\chapter{Resultados e Discussão}
\label{cap:resultados}

% -----------------------------------------------------------------
% Objetivo do capítulo:
% Apresentar os resultados obtidos com a ferramenta em dados
% reais e discutir o que eles dizem sobre a proposição central.
% O capítulo deve responder, de forma explícita:
%   - O que cada módulo, isolado, entrega?
%   - O que a abordagem híbrida agrega?
%   - O resultado é coerente com o conhecimento de
%     especialistas?
%   - Onde a ferramenta acerta, onde erra e por quê?
%
% Este capítulo é o último a ser preenchido. Por ora ficam
% apenas as seções com a função de cada uma anotada.
% -----------------------------------------------------------------


\section{Configuração experimental}
\label{sec:config_experimental}

% Conteúdo a desenvolver:
%   - Recorte da base usada nos experimentos.
%   - Período analisado, máquinas, lotes e fornecedores
%     considerados.
%   - Hiperparâmetros relevantes e versões dos scripts.


\section{Comportamento do módulo de anomalia}
\label{sec:resultados_eq}

% Conteúdo a desenvolver:
%   - Distribuição dos escores de anomalia.
%   - Casos em que o módulo detectou desvio operacional.
%   - Cobertura do conjunto de 14 sensores e impacto das
%     combinações SEM_MODELO.


\section{Comportamento do módulo estatístico}
\label{sec:resultados_mp}

% Conteúdo a desenvolver:
%   - Lotes, materiais, fornecedores e máquinas com maior
%     risco relativo.
%   - Padrões de cruzamento entre máquinas que reforçam
%     suspeita de matéria-prima.


\section{Diagnóstico híbrido em casos reais}
\label{sec:resultados_hibrido}

% Conteúdo a desenvolver:
%   - Exemplos de diagnóstico em casos documentados.
%   - Caso IBRAME 482530 (devolução por bolhas): o que o
%     sistema reportou e o que indica sobre o falso negativo
%     atual.
%   - Casos em que o diagnóstico foi confirmado pelos
%     especialistas.


\section{Comparação entre abordagens}
\label{sec:comparacao}

% Conteúdo a desenvolver:
%   - Anomalia isolada x estatística histórica isolada x
%     abordagem híbrida.
%   - Métricas de comparação adotadas e justificadas.
%   - Ganho marginal da integração e cenários em que ele se
%     manifesta.


\section{Validação externa contra o relatório de devoluções}
\label{sec:validacao_devolucoes}

% Conteúdo a desenvolver:
%   - Uso do relatório de devoluções de 2025 como ground truth
%     supervisionado.
%   - Concordância e divergência entre diagnóstico do sistema
%     e devolução registrada.
%   - Discussão sobre o que cada divergência indica.


\section{Discussão geral, ganhos e limitações}
\label{sec:discussao}

% Conteúdo a desenvolver:
%   - Em que o sistema avança em relação a abordagens
%     isoladas.
%   - Limites observados (cobertura de sensor, dependência da
%     qualidade dos dados, casos em que o diagnóstico é
%     inconclusivo).
%   - Implicações práticas para manutenção, qualidade e
%     engenharia.
